\scanpage{106}
\mspage{21}
% archive sheet number: 105
% source lines 1--2
et du calcul 4.2 des $\underline H^q_Y(F)$, impliquant que $E^{pq}_2=0$ si $q\ne 2d$, d'où $H^i_Y(X,F)\simeq E^{i-2d,2d}_2$.

% source lines 3--4
Faisant \deleted{\unsure{$F$}} $F=\mu_n^{\otimes d}$ et $i=2d$ dans (4.3.1), on trouve un homomorphisme

\[
H^{2d}_Y(X,\deleted{\unsure{\mu_n^{\otimes d}}}\Lambda(d))\simeq H^0(Y,\Lambda)\longrightarrow H^{2d}(X,\deleted{\unsure{\mu_n^{\otimes d}}}\Lambda(d))
\]

% source lines 6--11
et l'image de la section \deleted{1 \unsure{dans le premier}} \inserted{\unsure{constante} 1 de $\Lambda_Y$ sur $Y$} nous donne un élément canonique (4.3.2). Les assertions \deleted{\unsure{les}} de compatibilité faites \inserted{dans 4.3} sur cet élément \deleted{dans} se prouvent par des arguments standard, pour lesquels nous renvoyons à SGA 5 $\overline{\underline{\mathrm{IV}}}$.

% source lines 12--14
\uline{Remarque} 4.3.3. On a \inserted{\unsure{seulement}} utilisé l'hyp. de pureté faite pour \deleted{conclure} avoir \deleted{\unsure{les}} les \unsure{valeurs} de $\underline H^i_Y(\Lambda_X)$ \ldots\ .
